\changecolours{ScoutGreen}
\section[Referencias]{Referencias Bibliograficas}

% ------------------------------------------------------------------------------
\begin{frame}{Referências Bibliográficas Recomendadas}
  \scriptsize
  \itemseps[topsepi=0.5mm,itemsepi=1mm,labelsep=2mm,leftmargini=3mm]
  \begin{itemize}
    \item \textbf{ELMASRI, Ramez; NAVATHE, Shamkant B.} \textit{Sistemas de Banco de Dados}. 7. ed. São Paulo: Pearson Education do Brasil, 2018. \textbf{(Capítulo 9: Mapeamento do Modelo ER e EER para o Modelo Relacional -- Passos 5, 6, 7 e 8)}.
    \item \textbf{HEUSER, Carlos Alberto.} \textit{Projeto de Banco de Dados}. 6. ed. Porto Alegre: Bookman, 2009. \textbf{(Capítulo 4: Projeto Lógico -- Mapeamento de Associações $N:M$, Atributos Multivalorados e Generalizações)}.
    \item \textbf{SILBERSCHATZ, Abraham; KORTH, Henry F.; SUDARSHAN, S.} \textit{Sistema de Banco de Dados}. 7. ed. Rio de Janeiro: Elsevier, 2020. \textbf{(Capítulo 6: Redução para Esquemas Relacionais -- Conjuntos de Relacionamentos com Atributos Descritivos)}.
    \item \textbf{CODD, Edgar Frank.} \textit{A Relational Model of Data for Large Shared Data Banks}. Communications of the ACM (CACM), New York, v. 13, n. 6, p. 377--387, June 1970.
    \item \textbf{CHEN, Peter Pin-Shan.} \textit{The Entity-Relationship Model: Toward a Unified View of Data}. ACM Transactions on Database Systems (TODS), New York, v. 1, n. 1, p. 9--36, Mar. 1976.
  \end{itemize}
\end{frame}
